% ===== roadmap of this week =====
\begin{frame}{이번 주차 개요}
  \small
  \textbf{이 주차를 마치면}
  \begin{enumerate}
    \item 연속제어 환경에서 4주 파이프라인(환경 고르기 $\to$ 구현 $\to$ 평가
      $\to$ 발표)으로 팀 프로젝트를 설계·실행하고, ``시드 $\geq 3$ · 무작위
      정책 베이스라인 · 학습/평가 분리'' 프로토콜을 실제로 수행할 수 있다.
    \item ``학습 곡선이 올랐다''와 ``배운 정책이 무작위보다 낫다''가 서로 다른
      문장임을 구분하고, 학습 도중(탐험 포함) 성과와 학습 후 결정적($a=\mu$)
      평가의 격차를 11.2절의 분산 원리로 해석할 수 있다.
    \item 5항목 체크리스트와 1점(감상평) $\to$ 3점(반박 가능한 질문) 기준으로
      다른 팀의 발표를 심사하는 peer-review를 작성할 수 있다.
    \item 학기 전체를 ``모델 정의 $\to$ 틀린 정도를 수치화 $\to$ 개선'' 3단계로
      정리하고, 학기 이후 방향(오프라인 RL · 범용 로봇 정책 · sim-to-real)을
      자기 언어로 말할 수 있다.
  \end{enumerate}
  \vskip0.5em
  \begin{itemize}
    \item \kb{16.1} 팀 프로젝트: 발표 — Pendulum 실습, 시드 3개, 결정적 평가, 4주 파이프라인
    \item \kb{16.2} Peer-Review — 이번엔 네가 심사위원. ``우연히 좋은'' 패턴 5가지
    \item \kb{16.3} ML2 총정리 — 하나의 MDP를 네 경로로, 20문항 자가진단, 최신 동향
  \end{itemize}
  \vskip0.4em
  \kq{Week 9--15로 새 도구(신경망·시뮬레이션·트리 탐색)를 다 갖췄다.
  이 장부터는 배울 개념이 남지 않는다 — 도구를 로봇 시뮬레이션에 직접
  적용해, 그 결과를 정직하게 보고하는 것으로 학기를 마무리한다.}
\end{frame}
